\documentclass[11pt]{article}

\usepackage[margin=1in]{geometry}
\usepackage{amsmath,amssymb,booktabs,tabularx}

\title{Risk-neutral characteristic function}
\author{}
\date{}

\begin{document}
\maketitle

\noindent\textbf{Status.} Draft\\
\textbf{Plan item.} Risk-neutral characteristic function \(\phi_Q(u)\)\\
\textbf{Depends on.} \texttt{08\_pricing\_density\_normalization.tex},
\texttt{09\_martingale\_tilt\_equation.tex}\\
\textbf{Feeds.} Terminal-affine spectral option pricing, pricing diagnostics, put-call parity checks

\section{Purpose}

Derive the characteristic function of terminal log price under the normalized
terminal pricing measure.
This is the transform consumed by the spectral vanilla option pricer.

\section{Setup}

Let
\[
    \log S_T
    =
    \log S_t+d_{t,T}+X_{t,T}.
\]
Equivalently, if the deterministic terminal log price conditional on zero
future stochastic return is
\[
    \ell_{t,T}^{0}=\log S_t+d_{t,T},
\]
then
\[
    \log S_T=\ell_{t,T}^{0}+X_{t,T}.
\]
The pricing density is
\[
    M_T(\theta,\pi_+,\pi_-)
    =
    \exp\!\left(
        \theta X_{t,T}
        +\log(\pi_+)N_{t,T}^+
        +\log(\pi_-)N_{t,T}^-
        -K_{t,T}(\theta,\log\pi_+,\log\pi_-)
    \right),
\]
where \(\theta\) solves the martingale equation in
\texttt{09\_martingale\_tilt\_equation.tex}.

\section{Notation}

\begin{tabularx}{\textwidth}{lXX}
\toprule
Symbol & Meaning & Units or domain\\
\midrule
\(u\) & Fourier frequency for terminal log price & Real\\
\(\phi_Q(u)\) & Risk-neutral characteristic function of \(\log S_T\) & Complex, \(\phi_Q(0)=1\)\\
\(\theta\) & Solved martingale tilt & Real cumulant domain\\
\(\pi_\pm\) & Activity-premium multipliers & Strictly positive\\
\(K_{t,T}\) & Physical stochastic terminal cumulant & Finite at \(\theta+iu\) and \(\theta\)\\
\(\ell_{t,T}^{0}\) & Deterministic terminal log-price component & Real finite\\
\bottomrule
\end{tabularx}

\section{Assumptions used}

\begin{tabularx}{\textwidth}{XXX}
\toprule
Assumption & Where it is used & Why it is needed\\
\midrule
The terminal pricing density is normalized.
& Dividing by \(K_{t,T}(\theta,\log\pi_+,\log\pi_-)\).
& The characteristic function must be an expectation under a probability measure.\\
The martingale tilt has already been solved.
& Treating \(\theta\) as fixed in \(\phi_Q\).
& The risk-neutral transform depends on the no-arbitrage tilt.\\
The terminal log price splits into deterministic and stochastic components.
& Factoring out \(e^{iu\ell_{t,T}^{0}}\).
& Deterministic terminal log price contributes only a phase.\\
The cumulant is finite at \(\theta+iu\) and \(\theta\).
& Evaluating numerator and denominator transforms.
& The characteristic function must be finite.\\
\bottomrule
\end{tabularx}

\section{Result}

\noindent\textbf{Identity.} The risk-neutral characteristic function of
terminal log price is
\[
    \phi_Q(u)
    =
    \exp\!\left(
        iu\,\ell_{t,T}^{0}
        +K_{t,T}(\theta+iu,\log\pi_+,\log\pi_-)
        -K_{t,T}(\theta,\log\pi_+,\log\pi_-)
    \right).
\]
For terminal log return \(R_{t,T}=d_{t,T}+X_{t,T}\), replace
\(\ell_{t,T}^{0}\) by \(d_{t,T}\).

\section{Derivation}

By definition,
\[
    \phi_Q(u)
    =
    \mathbb E^Q\!\left[e^{iu\log S_T}\mid\mathcal F_t\right].
\]
Using \(\log S_T=\ell_{t,T}^{0}+X_{t,T}\),
\[
    \phi_Q(u)
    =
    e^{iu\ell_{t,T}^{0}}
    \mathbb E^Q\!\left[e^{iuX_{t,T}}\mid\mathcal F_t\right].
\]
Change measure with \(M_T\):
\[
    \mathbb E^Q\!\left[e^{iuX_{t,T}}\mid\mathcal F_t\right]
    =
    \mathbb E^P\!\left[
        e^{iuX_{t,T}}M_T(\theta,\pi_+,\pi_-)
        \mid\mathcal F_t
    \right].
\]
Substituting the density gives
\[
    \mathbb E^Q\!\left[e^{iuX_{t,T}}\mid\mathcal F_t\right]
    =
    \frac{
        \mathbb E^P\!\left[
            \exp\!\left(
                (\theta+iu)X_{t,T}
                +\log(\pi_+)N_{t,T}^+
                +\log(\pi_-)N_{t,T}^-
            \right)
            \mid\mathcal F_t
        \right]
    }{
        \exp\!\left(K_{t,T}(\theta,\log\pi_+,\log\pi_-)\right)
    }.
\]
The numerator is
\[
    \exp\!\left(K_{t,T}(\theta+iu,\log\pi_+,\log\pi_-)\right).
\]
Therefore
\[
\mathbb E^Q\!\left[e^{iuX_{t,T}}\mid\mathcal F_t\right]
=
    \exp\!\left(
        K_{t,T}(\theta+iu,\log\pi_+,\log\pi_-)
        -K_{t,T}(\theta,\log\pi_+,\log\pi_-)
    \right),
\]
and multiplying by the deterministic phase gives the stated formula.

At \(u=0\),
\[
    \phi_Q(0)
    =
    \exp\!\left(
        K_{t,T}(\theta,\log\pi_+,\log\pi_-)
        -K_{t,T}(\theta,\log\pi_+,\log\pi_-)
    \right)
    =1,
\]
so the transform is normalized.

\section{Domain and validity}

The formula is valid for real frequencies \(u\) such that the terminal
cumulant is finite at both \(\theta+iu\) and \(\theta\).
The solved \(\theta\) must satisfy the martingale equation, the activity premia
must be strictly positive, and the affine schedule domain must hold along every
interval.
The expression is the terminal-log-price characteristic function.
For terminal log returns, the deterministic phase uses \(d_{t,T}\) instead of
\(\ell_{t,T}^{0}\).

\section{Consequences}

\begin{itemize}
    \item Spectral option pricing can use a single normalized terminal characteristic function for calls and puts.
    \item The density denominator enforces \(\phi_Q(0)=1\).
    \item The \(iu\) shift is the Fourier analogue of the \(+1\) shift in the martingale equation.
    \item The deterministic terminal log-price path contributes only a phase factor.
\end{itemize}

\section{Checks}

\begin{tabularx}{\textwidth}{XXX}
\toprule
Check & Expected result & Evidence location\\
\midrule
Frequency argument
& The numerator cumulant is evaluated at \(\theta+iu\).
& \texttt{TerminalAffineCumulantArgument::from\_theta\_and\_frequency} in \texttt{crates/janiform-engine-pricing/src/spectral/pricer/terminal\_affine/cumulant.rs}\\
Characteristic assembly
& The characteristic value is the deterministic log-price phase times \(\exp(K(\theta+iu)-K(\theta))\).
& \texttt{solve\_terminal\_affine\_frequency\_output} in \texttt{crates/janiform-engine-pricing/src/spectral/pricer/terminal\_affine/core.rs}\\
Zero-frequency normalization
& \(\phi_Q(0)=1\) within numerical tolerance.
& Test \texttt{terminal\_affine\_measure\_change\_normalizes\_characteristic\_and\_martingale}\\
Long-variance dependence
& Changing the initial long-variance state changes the terminal-affine characteristic function.
& Test \texttt{long\_variance\_state\_changes\_terminal\_affine\_characteristic\_function}\\
\bottomrule
\end{tabularx}

\end{document}
